\documentstyle{article}
\begin{document}

\centerline{\bf A Method of Constructing the Real Canonical Form}
\centerline{\bf of a Hamiltonian Matrix with two Imaginary Eigenvalues}

\vskip 0.3cm

\centerline{Rodney Coleman}
\centerline{Laboratoire LMC-IMAG}
\centerline{46 Avenue F\'{e}lix Viallet}
\centerline{38031 Grenoble Cedex, France}
\centerline{Email: coleman@imag.fr}

\vskip 0.3cm

Hamiltonian matrices play an important role in Hamiltonian systems of differential equations. Such matrices are of even order and can be written $A = JS$,
where $S$ is symmetric and
$$
J = J_{2n} = \left\lgroup \begin{array}{cc}
		0 & I_n\\
		-I_n & 0
		\end{array} \right\rgroup
$$

\noindent It is not in general true that conjugation by a non-singular matrix
$T$ produces a Hamiltonian matrix; however this is the case if $T$ is
symplectic, i.e.,
$$
T^tJT=J
$$
\noindent It is natural to ask whether, for a given Hamiltonian matrix
$A$, we can find a symplectic matrix $T$ such that $\tilde{A} =
T^{-1}AT$ is of a particularly simple form and, if so, whether this
form is unique in some sense. The answer to this question is
affirmative. Thus in the same way as we speak of the real Jordan
canonical form of a general matrix we can speak of the real canonical
form of a Hamiltonian matrix.\\

\noindent A particular case is that of a Hamiltonian matrix with a pair
of imaginary eigenvalues $\pm  i \beta $. Unlike certain cases, for
example where the matrix has a pair of real eigenvalues $\pm \alpha $,
here the construction of a canonical form is not at all easy. The
method which we propose differs fundamentally in approach from that of
N. Burgoyne and R. Cushman.  It is based on some ideas of A.Laub and
K.~R. Meyer.

\end{document}
